\documentclass[12pt,a4paper]{report}
% ============================= DOCUMENT CLASS & GEOMETRY =============================
\usepackage[left=1.5in,right=1.0in,top=1.0in,bottom=1.0in]{geometry}
\usepackage{times}
\usepackage{setspace}
\renewcommand{\baselinestretch}{1.4}

% ============================= ENHANCED PACKAGES =============================
% Typography & Formatting
\usepackage[titles]{tocloft}
\usepackage{afterpage}
\usepackage{sectsty}
\usepackage{titlesec}
\usepackage{lettrine}
\usepackage{lipsum}

% Graphics & Figures
\usepackage{graphicx}
\usepackage{xcolor}
\usepackage[export]{adjustbox}
\usepackage{rotating}
\usepackage{float}

% Mathematics & Science
\usepackage{amsfonts}
\usepackage{amstext}
\usepackage{amssymb}
\usepackage{amsmath,bm}
\usepackage{braket}
\usepackage{cases}

% Algorithms & Code
\usepackage{algorithm}
\usepackage{algpseudocode}

% Tables & Lists
\usepackage{booktabs}
\usepackage{longtable}
\usepackage{lscape}

% Language & References
\usepackage[english]{babel}
\usepackage{parskip}
\usepackage{url}
\usepackage[acronym]{glossaries}
\usepackage{nomencl}

% Captions & Typography
\usepackage[font=small,format=plain,labelfont=bf,up,textfont=it,up]{caption}
\usepackage[T1]{fontenc}

% Enhanced Hyperlinks
\usepackage[colorlinks=true, allcolors=blue, bookmarks=true, bookmarksnumbered=true]{hyperref}

% ============================= CUSTOM FORMATTING =============================
% Table of Contents Depth
\setcounter{tocdepth}{3}
\setcounter{secnumdepth}{3}

% Section Font Formatting
\allsectionsfont{\usefont{OT1}{times}{bc}{n}\bf\selectfont}

% Title Formatting - Main Chapters
\titleformat{\chapter}[display]{\vspace{7.5cm}\Large\bfseries\centering}{\chaptername~\thechapter}{1ex}{}

% Enhanced TOC Formatting
\renewcommand{\contentsname}{\textbf{TABLE OF CONTENTS\normalsize\centering}}
\renewcommand{\listfigurename}{LIST OF FIGURES}
\renewcommand{\listtablename}{LIST OF TABLES}

% Paragraph Spacing
\parskip 1ex

% ============================= BEGIN DOCUMENT =============================
\begin{document}

% ============================= TITLE PAGE (Page 1) =============================
\begin{titlepage}
	\centering

	% Title
	{\LARGE \textbf{ Quantum Steering } \par}

	\vspace{1.5cm}

	% University Logo
	\includegraphics[width=7.7cm]{QAU_enhanced_Logo} % Update filename as needed

	\vspace{1.2cm}

	% Author
	{\large \textbf{BY} \par}
	{\Large \textbf{[Hafsa Zia]} \par}
	\vspace{0.6cm}

	% Supervisor
	{\large \textbf{Supervisor:} \\}
	{\Large \textbf{Dr. Shakir Ullah} \par}

	\vspace{4cm}

	% Department and University
	{\Large\textbf{Department of Physics} \par}
	{\textbf{Quaid-i-Azam University, Islamabad, Pakistan} \par}
	{(Spring 2025--2026) \par}

\end{titlepage}

% ============================= TITLE PAGE (Page 2 - Formal) =============================
\newpage
\begin{titlepage}
	\centering

	% Title
	{\LARGE \textbf{Continuous-Variable Quantum Steering in a Raman-Driven Quantum Beat Laser} \par}
	\vspace{0.8cm}

	\includegraphics[width=7.5cm]{QAU_enhanced_Logo}
	\vspace{0.5cm}

	% Author
	{\large \textbf{BY} \par}
	{\Large \textbf{[Your Name]} \\}
	{\centering\large (Reg. \# [Your Registration Number])}
	\vspace{0.5cm}

	% Supervisor
	{\large \textbf{Supervisor:} \\}
	{\Large \textbf{Dr. Shakir Ullah} \par}

	\vspace{0.8cm}
		% Dissertation Note
	\begin{spacing}{1.2}
		{\textbf{\textit{A dissertation submitted in partial fulfillment of the requirements for the degree of}}}
		{\Large \textit{\textbf{Master of Science in Physics}} \par}
		{\textbf{\textit{at the Quaid-i-Azam University, Islamabad, Pakistan}} \par}
		{\textbf{\textit{2025--2026}} \par}
	\end{spacing}
		\vspace{1cm}
		% Department and University
	{\Large\textbf{Department of Physics} \\}
	{\textbf{Quaid-i-Azam University, Islamabad, Pakistan} \par}

\end{titlepage}

% ============================= CERTIFICATION PAGE (Page 3) =============================
\newpage
\begin{titlepage}
	 \begin{center}
		{\Large \textbf{Research Completion Certificate}}
	\end{center}
	\vspace{0.4cm}
	This is to certify that \textbf{[Hafsa Zia]} (Reg. \# [Your Registration Number]) has carried out the work contained in this dissertation under my supervision and is accepted by the Department of Physics, Quaid-i-Azam University, Islamabad as satisfying the dissertation requirement for the degree of Master of Science in Physics.\\
	\vspace{1cm}

	% Right-aligned block (Supervisor info)
	\begin{flushright}
		\textbf{Supervisor:}\\[6pt]
		\vspace{1cm}
		\rule{3.5cm}{0.4pt}\\
		\textbf{Dr. Shakir Ullah}\\
		Department of Physics\\
		Quaid-i-Azam University,\\
		Islamabad, Pakistan.
	\end{flushright}
	\vspace{1cm}
	% Left-aligned block (Chairman info)
	\textbf{Submitted through:}\\[12pt]

	\begin{flushleft}
		\rule{3.5cm}{0.4pt}\\
		\textbf{Prof. Dr. [Chairman Name]}\\
		\textbf{Chairman}\\
		Department of Physics\\
		Quaid-i-Azam University,\\
		Islamabad, Pakistan.
	\end{flushleft}
\end{titlepage}

% ============================= DEDICATION PAGE (Page 4) =============================
\newpage
\begin{titlepage}
 \begin{center}
	{\Large \textbf{Dedication}}
\end{center}
\vspace{0.4cm}
\centering {I dedicate this work to my parents, whose unconditional support, encouragement, and sacrifices have been the foundation of my academic journey. Their belief in the pursuit of knowledge and excellence continues to inspire me.
}.\\
\end{titlepage}

% ============================= FRONT MATTER - ROMAN NUMERALS =============================
\pagenumbering{roman}
\setcounter{page}{8}
\newpage

% ============================= ACKNOWLEDGMENTS =============================
\addcontentsline{toc}{chapter}{Acknowledgments}
\chapter*{Acknowledgments}
I express my sincere gratitude to my supervisor, Dr. Shakir Ullah, for his invaluable guidance, insightful discussions, and continuous encouragement throughout this research. His expertise in quantum optics and quantum information processing has been instrumental in shaping this work.

I would like to thank the Department of Physics at Quaid-i-Azam University for providing the necessary research facilities and computational resources. Special thanks to my colleagues and research group members for their collaborative spirit and constructive feedback.

I am grateful to the Higher Education Commission (HEC), Pakistan, for their financial support through research grants. Finally, I acknowledge the contributions of previous researchers whose foundational work in quantum steering and laser physics has provided the theoretical basis for this investigation.

\newpage

% ============================= ABSTRACT =============================
\addcontentsline{toc}{chapter}{Abstract}
\chapter*{Abstract}

\noindent \textbf{Title:} Continuous-Variable Quantum Steering in a Raman-Driven Quantum Beat Laser\\[0.5cm]

Quantum steering—an intrinsically asymmetric quantum correlation intermediate between entanglement and Bell nonlocality—represents a fundamental resource for quantum information processing, particularly in one-sided device-independent quantum key distribution and asymmetric quantum networks. This dissertation investigates the theoretical framework for generating, controlling, and characterizing continuous-variable (CV) quantum steering in a Raman-driven quantum beat laser (RDQBL) system.

We develop a comprehensive theoretical model of the RDQBL operating in the regime of strong driving fields and weak cavity damping, where two external classical fields interact with a three-level atomic system in a Raman configuration to produce correlated photon pairs in two cavity modes. Through analytical and numerical approaches, we derive the master equation governing the system dynamics and compute the density matrix for the output cavity field using both interaction picture techniques and covariance matrix formalism.

Key findings demonstrate that:
\begin{itemize}
	\item Continuous-variable quantum steering is periodically generated between the two cavity modes as they propagate, with tunability through system parameters.
	\item The directionality of steering (one mode steering another but not vice versa) can be controlled via the relative phase between classical coupling fields and the Rabi frequency.
	\item Cavity damping introduces decoherence but the hierarchical structure of quantum correlations (Discord $\supseteq$ Entanglement $\supseteq$ Steering $\supseteq$ Bell Nonlocality) persists in realistic conditions.
	\item Non-classical properties of initial cavity modes and purity effects substantially influence steering dynamics and robustness.
\end{itemize}

This work bridges the gap between quantum optics—focusing on coherent light-matter interactions—and quantum information theory by providing the theoretical foundation and quantitative measures for exploiting steering resources in active photonic systems. The tunability of steering in RDQBL establishes it as a promising platform for next-generation quantum communication technologies and quantum metrology applications.

\textbf{Keywords:} Quantum steering, continuous variables, Raman-driven quantum beat laser, cavity QED, quantum correlations hierarchy, quantum communication.

\newpage

% ============================= ABBREVIATIONS =============================
\addcontentsline{toc}{chapter}{List of Abbreviations}
\chapter*{List of Abbreviations}

\begin{tabular}{ll}
CV & Continuous Variable \\
DV & Discrete Variable \\
RDQBL & Raman-Driven Quantum Beat Laser \\
QBL & Quantum Beat Laser \\
CEL & Correlated Emission Laser \\
QIT & Quantum Information Theory \\
QED & Quantum Electrodynamics \\
CQED & Cavity Quantum Electrodynamics \\
EPR & Einstein-Podolsky-Rosen \\
1sDI & One-Sided Device-Independent \\
QKD & Quantum Key Distribution \\
1sDI-QKD & One-Sided Device-Independent Quantum Key Distribution \\
PPT & Positive Partial Transpose \\
CV & Continuous Variable \\
TMGS & Two-Mode Gaussian State \\
NOON & N-Photon State \\
HOM & Hong-Ou-Mandel (interference) \\
PMF & Probability Mass Function \\
PDF & Probability Density Function \\
FWHM & Full Width at Half Maximum \\
SNR & Signal-to-Noise Ratio \\
\end{tabular}

\newpage

% ============================= TABLE OF CONTENTS =============================
\renewcommand{\contentsname}{\textbf{TABLE OF CONTENTS}}
\tableofcontents

\titleformat{\chapter}[display]{\large\bfseries\raggedright}{\chaptername~\thechapter}{1ex}{}[{\titlerule[1pt]}]

\newpage

% ============================= LIST OF FIGURES =============================
\addcontentsline{toc}{chapter}{List of Figures}
\listoffigures

\newpage

% ============================= LIST OF TABLES =============================
\addcontentsline{toc}{chapter}{List of Tables}
\listoftables

\newpage

% ============================= MAIN CONTENT - ARABIC NUMERALS =============================
\pagenumbering{arabic}
\titleformat{\chapter}[display]{\vspace{7.5cm}\Large\bfseries\centering}{\chaptername~\thechapter}{1ex}{}

% ============================= CHAPTER 1: INTRODUCTION =============================
%=============================================================================================
%                                  Chapter 1: Introduction
%=============================================================================================

\chapter{Introduction}
\label{chp1}
\newpage

\section{Background and Motivation}

Quantum correlations form the foundational basis of quantum information science and are essential to the operation of emerging quantum technologies such as quantum computing, quantum communication, and quantum cryptography. Among the various forms of quantum correlations, \textbf{quantum steering} stands out as a particularly intriguing and asymmetric nonclassical effect that occupies a unique position between entanglement and Bell nonlocality. Unlike entanglement and Bell nonlocality, which are symmetric under exchange of the two parties, quantum steering is inherently directional—one party may be able to steer another, but the reverse might not hold, highlighting an asymmetry with profound implications for quantum information protocols.

The concept of quantum steering was first introduced by Erwin Schrödinger in 1935 as a direct response to the Einstein-Podolsky-Rosen (EPR) paradox \cite{SchrodingerE35}, which questioned the completeness of quantum mechanics. Schrödinger originally termed this phenomenon ``steering'' to describe the scenario where one spatially distant system can influence the state of another through local measurements, despite their physical separation. This directional nature of steering has remained largely unexplored compared to its symmetric counterparts until recent years, when its potential applications in quantum cryptography became apparent.

In the context of continuous-variable (CV) quantum systems, steering provides a powerful resource for quantum information processing. Unlike discrete-variable (DV) systems that rely on photon number states and post-selection techniques, CV systems allow for the generation of entangled states deterministically through optical parametric processes and laser-based systems. The measurable quantities in CV systems—such as position, momentum, or the quadrature components of light fields—can be continuously varied, enabling more efficient and robust quantum protocols.

\subsection{Quantum Steering: Definition and Significance}

Quantum steering can be formally defined as follows: A bipartite quantum system in state $\rho_{AB}$ exhibits steering of subsystem B by subsystem A if Alice, through her local measurements on subsystem A, can influence the conditional states of Bob's subsystem B in a manner that cannot be explained by any local hidden-variable model. This asymmetric definition stands in sharp contrast to entanglement (which is symmetric) and Bell nonlocality (which is symmetric under relabeling).

The practical significance of steering manifests in several quantum information protocols:

\begin{itemize}
	\item \textbf{One-Sided Device-Independent (1sDI) Quantum Key Distribution (QKD):} In traditional QKD schemes, both Alice and Bob must trust their measurement devices. However, in 1sDI-QKD protocols, security can be ensured even when only one party operates a trusted measurement device \cite{BranciardC12}. This substantially reduces the hardware requirements and vulnerability to side-channel attacks.

	\item \textbf{Asymmetric Quantum Networks:} In quantum networks where nodes have unequal capabilities or trust levels, directional steering enables secure communication with minimal trusted resources.

	\item \textbf{Quantum Metrology:} The asymmetric nature of steering can be exploited for parameter estimation in scenarios where only one system can be precisely manipulated.
\end{itemize}

\subsection{Continuous Variables vs. Discrete Variables}

Continuous-variable quantum systems offer several practical advantages over discrete-variable systems:

\begin{enumerate}
	\item \textbf{Deterministic Generation:} Entangled CV states can be generated deterministically (e.g., through squeezed light and parametric down-conversion), without requiring probabilistic post-selection as in DV systems.

	\item \textbf{Optical Compatibility:} CV states are naturally compatible with existing optical infrastructure, including fiber-optic telecommunications networks, enabling long-distance quantum communication.

	\item \textbf{Homodyne Detection:} Measurement of CV observables can be performed using well-established homodyne detection techniques, which are simpler and more efficient than single-photon counting in DV systems.

	\item \textbf{Scalability:} The continuous nature of CV variables allows for greater scalability in quantum computing and information processing implementations.
\end{enumerate}

\section{Active vs. Passive Quantum Systems}

Recent progress in quantum information science has highlighted the importance of distinguishing between passive and active quantum systems:

\textbf{Passive Systems:} Linear optical elements such as beam splitters, interferometers, and directional couplers. While useful for manipulating quantum states, passive systems are limited in their ability to generate new quantum correlations. Their primary utility lies in redistributing or transforming existing correlations.

\textbf{Active Systems:} Laser-based systems and quantum emitters that actively generate quantum correlations through coherent light-matter interactions. Examples include:
\begin{itemize}
	\item Quantum Beat Lasers (QBL) - where two classical driving fields create quantum beat phenomena
	\item Correlated Emission Lasers (CEL) - three-level atomic systems producing entangled photon pairs
	\item Parametric Oscillators - producing squeezed and entangled light through nonlinear processes
\end{itemize}

Active systems are superior for generating robust, tunable quantum correlations that can persist despite dissipation and decoherence.

\section{Raman-Driven Quantum Beat Laser System}

The Raman-driven quantum beat laser (RDQBL) represents an advanced active quantum system that combines several beneficial features for steering generation:

\subsection{Raman Process Fundamentals}

The Raman process involves coherent population transfer between two long-lived atomic states via an intermediate excited state, typically mediated by two classical laser fields. Key advantages of Raman processes include:

\begin{itemize}
	\item \textbf{Minimal Population in Excited State:} Indirect coupling through the Raman process bypasses significant population accumulation in the excited state, thereby reducing spontaneous emission losses and decoherence.

	\item \textbf{Enhanced Coherence:} By avoiding excited state populations, Raman systems maintain superior coherence properties essential for quantum correlations.

	\item \textbf{Optical Gain:} Raman processes facilitate stimulated emission that can lead to optical amplification in specific modes, creating a gain-driven environment favorable for quantum state engineering.
\end{itemize}

\subsection{Quantum Beat Phenomena}

When a quantum system is driven by two classical fields with frequencies $\nu_1$ and $\nu_2$, the beat frequency $\Omega_{\text{beat}} = |\nu_1 - \nu_2|$ emerges in the dynamical evolution. Quantum beats manifest as periodic oscillations in observables and correlation functions, allowing for time-dependent control of quantum properties.

\section{Problem Statement and Research Objectives}

\subsection{Motivation for This Work}

While significant theoretical progress has been made in understanding quantum correlations in passive linear systems (e.g., coupled waveguides) and in understanding entanglement in active laser systems (CEL, QBL), the specific problem of \textbf{asymmetric quantum steering in Raman-driven quantum beat lasers remains largely unexplored}.

Most prior work has focused on:
\begin{itemize}
	\item Entanglement generation in waveguides (without addressing asymmetry)
	\item Bell nonlocality in laser systems
	\item General quantum discord without directional considerations
\end{itemize}

In contrast, this work addresses the critical gap: How can one systematically generate, control, and quantify \textbf{asymmetric quantum steering} in an active Raman-driven system?

\subsection{Main Research Objectives}

The primary objectives of this research are:

\begin{enumerate}
	\item \textbf{Develop a comprehensive theoretical framework} for continuous-variable quantum steering in the RDQBL system, including:
	\begin{itemize}
		\item Detailed Hamiltonian formulation for the four-level atomic model
		\item Master equation derivation including cavity damping and atomic decay
		\item Analytical and numerical solutions for system dynamics
	\end{itemize}

	\item \textbf{Quantify quantum steering} using established criteria:
	\begin{itemize}
		\item Covariance matrix analysis for Gaussian states
		\item Logarithmic negativity measures
		\item Verification of the hierarchical structure: Discord $\supseteq$ Entanglement $\supseteq$ Steering $\supseteq$ Bell Nonlocality
	\end{itemize}

	\item \textbf{Investigate steering tunability} through system parameters:
	\begin{itemize}
		\item Effects of Rabi frequency of classical coupling fields
		\item Relative phase control between driving fields
		\item Cavity damping rates and their decoherence effects
		\item Non-classical and purity properties of initial states
	\end{itemize}

	\item \textbf{Demonstrate asymmetric steering directivity}:
	\begin{itemize}
		\item Show that Mode 1 can steer Mode 2 but not vice versa
		\item Achieve selective steering through parameter optimization
		\item Map parameter regimes for maximal steering asymmetry
	\end{itemize}

	\item \textbf{Assess robustness under realistic conditions}:
	\begin{itemize}
		\item Evaluate steering persistence under cavity losses
		\item Study effects of atomic decay rates
		\item Determine critical thresholds for steering survival
	\end{itemize}
\end{enumerate}

\section{Significance and Expected Contributions}

This research is expected to contribute to quantum information science in several ways:

\begin{itemize}
	\item \textbf{Theoretical Foundation:} Provides a complete theoretical treatment of steering in active laser systems, filling a gap between quantum optics and quantum information theory.

	\item \textbf{Practical Platform Identification:} Demonstrates RDQBL as a viable, experimentally-realizable platform for generating tunable, robust quantum steering resources.

	\item \textbf{Asymmetric Protocol Design:} Offers insights into exploiting directionality for designing advanced 1sDI-QKD and asymmetric quantum network protocols.

	\item \textbf{Decoherence Resilience:} Shows how gain and coherence control can mitigate decoherence—a critical concern for practical quantum technologies.

	\item \textbf{Interdisciplinary Bridge:} Connects quantum optics, laser physics, and quantum information science, enabling future interdisciplinary research.
\end{itemize}

\section{Thesis Layout}

This dissertation is organized as follows:

\begin{itemize}
	\item \textbf{Chapter \ref{chp2}:} Presents comprehensive background on quantum mechanics fundamentals, continuous-variable systems, Gaussian state formalism, quantum steering theory, and the RDQBL system architecture. Both theoretical foundations and practical aspects are covered.

	\item \textbf{Chapter \ref{chp3}:} Introduces the physical RDQBL model with its Hamiltonian, four-level atomic configuration, and master equation. Derivations of covariance matrices and steering quantification measures are presented. Both analytical and numerical approaches for solving the dynamics are discussed.

	\item \textbf{Chapter \ref{chp4}:} Presents comprehensive numerical results showing how quantum steering evolves in time for various parameter regimes. Effects of cavity damping, Rabi frequencies, initial state properties, and relative phases are systematically investigated and visualized.

	\item \textbf{Chapter \ref{chp5}:} Provides in-depth discussion of results, comparison with related work, physical interpretations of steering dynamics, and implications for quantum technologies.

	\item \textbf{Chapter \ref{chp6}:} Summarizes the main findings and presents conclusions regarding the viability of RDQBL as a steering source. Future research directions and technological applications are outlined.
\end{itemize}

\section{Key Contributions of This Work}

To summarize, the key original contributions of this dissertation are:

\begin{enumerate}
	\item First comprehensive treatment of \textbf{continuous-variable quantum steering} (not just entanglement) in a Raman-driven system
	\item Demonstration of \textbf{tunable, asymmetric steering} through accessible laser parameters
	\item Rigorous \textbf{hierarchical analysis} of quantum correlations in RDQBL
	\item Quantitative assessment of \textbf{steering robustness} under realistic dissipation
	\item Practical roadmap for \textbf{experimental implementation} and quantum technology applications
\end{enumerate}

This work thus aspires to deepen our understanding of quantum steering as a resource for quantum information processing and to establish Raman-driven systems as powerful platforms for realizing steering-based quantum technologies in realistic, noisy environments.



% ============================= CHAPTER 2: LITERATURE REVIEW & THEORY =============================
%=============================================================================================
%                           Chapter 2: Theory and Literature Review
%=============================================================================================

\chapter{Quantum Correlations, Continuous Variables, and Quantum Steering}
\label{chp2}
\newpage

% ============================= SECTION 2.1: INTRODUCTION =============================
\section{Introduction}

Quantum correlations form the fundamental basis of quantum information science, representing non-classical relationships between subsystems that have no classical analogues. Unlike classical information where correlations arise from shared randomness, quantum correlations emerge from the intrinsic non-commutativity of quantum observables and the superposition principle. This chapter establishes the theoretical foundation for understanding quantum steering in continuous-variable systems, with particular emphasis on the RDQBL platform.

We begin by reviewing the hierarchy of quantum correlations—a central organizational principle that categorizes different types of quantum correlations by their strength and scope. We then shift focus to continuous-variable systems and their advantages over discrete-variable approaches. Finally, we examine quantum steering specifically: its mathematical formulation, quantification criteria for Gaussian states, and its unique asymmetric properties that distinguish it from entanglement and Bell nonlocality.

% ============================= SECTION 2.2: LITERATURE REVIEW =============================
\section{Literature Review: Quantum Correlations in Quantum Information}

\subsection{Overview of Quantum Steering}

Quantum correlations have been at the heart of foundational quantum mechanics debates since Einstein, Podolsky, and Rosen (EPR) challenged the completeness of quantum theory in 1935 \cite{einsteinPodolskyRosen1935}. Their paradox highlighted the apparent non-locality inherent in quantum mechanics—that measurements on one system can instantaneously affect another distant system.

Over the past century, several types of quantum correlations have been identified and rigorously formalized:

\begin{enumerate}
	\item \textbf{Quantum Entanglement:} Two or more subsystems are in a non-separable state such that the overall state cannot be written as a product of individual subsystem states. Entanglement is symmetric—if subsystem A is entangled with B, then B is entangled with A.

	\item \textbf{Bell Nonlocality:} A stronger form of correlation where measurements on one system cannot be explained by any local hidden-variable model. Violates Bell inequalities. Symmetric under party exchange.

	\item \textbf{Quantum Steering:} An asymmetric correlation where one party's measurements can influence (``steer'') the conditional state of another distant party in a manner incompatible with local realism. First formally defined by Schrödinger \cite{schrodingerDiscussion1935} as a response to the EPR paradox.

	\item \textbf{Quantum Discord:} A broader measure of quantum correlations that captures non-classical correlations beyond entanglement. Present in many separable states. Most general measure.
\end{enumerate}

\subsection{Hierarchy of Quantum Correlations: Discord $\supseteq$ Entanglement $\supseteq$ Steering $\supseteq$ Bell Nonlocality}

A fundamental insight in quantum information theory is that different types of quantum correlations form a strict hierarchy \cite{brunnerBellNonlocality2014}:

\begin{equation}
\text{Discord} \supseteq \text{Entanglement} \supseteq \text{Steering} \supseteq \text{Bell Nonlocality}
\end{equation}

where $X \supseteq Y$ means X is a superset of Y (i.e., all states exhibiting Y also exhibit X, but not vice versa).

\textbf{Bell Nonlocality:} The most restrictive class. If a state violates Bell inequalities, it exhibits all forms of correlations below it.

\textbf{Steering:} Intermediate in strength. Asymmetric—one party may steer another without the converse being true. Every steerable state exhibits entanglement, but not every entangled state is steerable.

\textbf{Entanglement:} States that are non-separable. Symmetric under party exchange. Always exhibits discord, but may not exhibit steering or nonlocality.

\textbf{Discord:} Widest class, present in many separable states. Classical correlations can coexist with quantum discord. Some separable states exhibit discord.

This hierarchy has profound implications: focusing specifically on steering (our work) targets the optimal balance between correlation strength and operational utility, particularly for one-sided device-independent protocols where only one party needs trusted devices.

% ============================= SECTION 2.3: ENTANGLEMENT IN CV =============================
\section{Entanglement in Continuous Variables}

\subsection{Advantages of Continuous-Variable Systems}

Continuous-variable (CV) quantum information processing offers significant practical advantages over discrete-variable (DV) systems:

\begin{itemize}
	\item \textbf{Deterministic Generation:} Entangled CV states can be generated deterministically through parametric down-conversion and squeezed light, without requiring post-selection.

	\item \textbf{Optical Compatibility:} CV observables (quadrature amplitudes) are naturally measured via homodyne detection using standard optical components. Easy integration with existing fiber-optic infrastructure.

	\item \textbf{Scalability:} Continuous variables allow more quantum information encoding per physical system compared to qubits.

	\item \textbf{Experimental Simplicity:} Homodyne detection is experimentally simpler than single-photon counting required in DV systems.

	\item \textbf{Gaussian State Framework:} Many CV states and operations can be represented using Gaussian formalism, enabling efficient computation and analysis.
\end{itemize}

\subsection{Quadrature Operators and Gaussian States}

The fundamental observables in CV quantum optics are the \textbf{quadrature operators}:
\begin{equation}
X = \frac{1}{\sqrt{2}}(a + a^\dagger), \quad P = \frac{i}{\sqrt{2}}(a^\dagger - a)
\end{equation}

These satisfy the canonical commutation relation:
\begin{equation}
[X, P] = i
\end{equation}

\textbf{Gaussian States} are quantum states whose Wigner function (quasi-probability distribution) is Gaussian. For bipartite systems, Gaussian states are fully characterized by:
\begin{enumerate}
	\item Mean values (first moments): $\langle X_1 \rangle, \langle P_1 \rangle, \langle X_2 \rangle, \langle P_2 \rangle$
	\item Covariance matrix (second-order moments)
\end{enumerate}

The two-mode covariance matrix is:
\begin{equation}
V = \begin{pmatrix}
V_1 & C_{12} \\
C_{12}^T & V_2
\end{pmatrix}
\end{equation}

where each $V_j$ is a $2 \times 2$ block:
\begin{equation}
V_j = \begin{pmatrix}
\langle X_j^2 \rangle - \langle X_j \rangle^2 & \langle\{X_j, P_j\}\rangle/2 - \langle X_j \rangle\langle P_j \rangle \\
\langle\{X_j, P_j\}\rangle/2 - \langle X_j \rangle\langle P_j \rangle & \langle P_j^2 \rangle - \langle P_j \rangle^2
\end{pmatrix}
\end{equation}

\subsection{Quantifying Entanglement in CV: Logarithmic Negativity}

For Gaussian states, entanglement is quantified by \textbf{logarithmic negativity}:
\begin{equation}
E = \max\{0, -\log_2(2\tilde{\eta}_-)\}
\end{equation}

where $\tilde{\eta}_-$ is the smallest symplectic eigenvalue of the partial transpose of the covariance matrix.

A state is entangled if $E > 0$. The metric $\tilde{\eta}_-$ is computed from:
\begin{equation}
2\tilde{\eta}_{\pm}^2 = \zeta \pm \sqrt{\zeta^2 - 4\det V}
\end{equation}
where $\zeta = \det V_1 + \det V_2 - 2\det C_{12}$.

\subsection{Coherent and Squeezed States}

Important classes of CV states include:

\textbf{Coherent States:} Pure states minimizing uncertainty, $|\alpha\rangle = e^{-|\alpha|^2/2}\sum_{n=0}^{\infty} \frac{\alpha^n}{\sqrt{n!}}|n\rangle$.

\textbf{Squeezed States:} States with reduced variance in one quadrature at the expense of increased variance in the orthogonal quadrature. Squeeze parameter $r$ quantifies the squeezing degree.

These states form the foundation for entanglement generation in parametric processes and laser systems.

% ============================= SECTION 2.4: STEERING IN CONTINUOUS VARIABLES =============================
\section{Quantum Steering in Continuous Variables}

\subsection{Formal Definition and Steering Criteria for Gaussian States}

Quantum steering is formalized as follows: Party A (Alice) can steer Party B (Bob) if Alice's local measurements on her subsystem produce conditional states for Bob that cannot be explained by any local hidden-variable model compatible with realism.

For Gaussian states with covariance matrix $V$, the steering criterion derived from the PPT (Positive Partial Transpose) condition is \cite{kogiastQuantificationGaussianQuantum2015}:
\begin{equation}
S^{A \to B} = \max\left\{0, \frac{1}{2}\log_2\frac{\det V_A}{4\det V_{out}^T}\right\}
\end{equation}

where:
\begin{itemize}
	\item $V_A = \det V_1$ is the covariance determinant for Alice's subsystem
	\item $V_{out}^T$ is evaluated over the conditional states Bob can access
\end{itemize}

If $S^{A \to B} > 0$, steering from A to B is demonstrated.

\subsection{Asymmetry of Steering: The Key Distinction}

The defining characteristic of steering is its \textbf{asymmetry}. Unlike entanglement and Bell nonlocality (which are symmetric), steering can be directional:

\begin{equation}
S^{A \to B} > 0 \text{ while } S^{B \to A} = 0 \text{ (Asymmetric Steering)}
\end{equation}

This directionality has profound practical implications:
\begin{itemize}
	\item In asymmetric steering, only one party needs a trusted measurement device
	\item The steered party (Bob) can be untrusted or even adversarial
	\item Perfect for one-sided device-independent protocols
\end{itemize}

\subsection{Entanglement versus Steering}

While related, steering and entanglement are distinct:

\begin{itemize}
	\item \textbf{Entanglement:} Correlation structure of the state itself, symmetric, quantifies by logarithmic negativity
	\item \textbf{Steering:} Ability to influence conditional states through measurement, asymmetric, directional
\end{itemize}

An entangled state may or may not be steerable in both directions. A steerable state is necessarily entangled, but an entangled state need not be steerable (if both parties have local hidden-variable descriptions).

% ============================= SECTION 2.5: EFFECT OF LOSS AND DECOHERENCE =============================
\section{Effect of Loss and Decoherence on Quantum Correlations}

\subsection{Lindblad Master Equation and Loss Operators}

Real quantum systems are open systems interacting with their environment. The dynamics are governed by the Lindblad master equation:
\begin{equation}
\frac{d\rho}{dt} = -\frac{i}{\hbar}[H, \rho] + \sum_k \left(L_k \rho L_k^\dagger - \frac{1}{2}\{L_k^\dagger L_k, \rho\}\right)
\end{equation}

For a two-mode cavity system with loss rates $\kappa_1, \kappa_2$, the loss Liouvillian is:
\begin{equation}
\mathcal{L}_{\text{loss}}[\rho] = \kappa_1\left(a_1\rho a_1^\dagger - \frac{1}{2}\{a_1^\dagger a_1, \rho\}\right) + \kappa_2\left(a_2\rho a_2^\dagger - \frac{1}{2}\{a_2^\dagger a_2, \rho\}\right)
\end{equation}

\subsection{Robustness of Entanglement versus Steering in Waveguides}

Previous work on coupled lossy waveguides \cite{RaiDasAgarwal2010} demonstrated:

\begin{itemize}
	\item Entanglement shows considerable robustness against material loss
	\item Loss parameters $\gamma/J$ up to 1/10 still preserve entanglement
	\item Logarithmic negativity decays but does not vanish rapidly
	\item Different initial states (photon number, NOON, squeezed) show varying robustness
\end{itemize}

Key finding: \textbf{Waveguide structures are reasonably robust against loss effects and appropriate for quantum circuits}.

This provides motivation for investigating steering robustness in similar systems—if entanglement persists, can steering (a potentially stronger resource) also be preserved?

% ============================= SECTION 2.6: OPTICAL WAVEGUIDES AS QM SYSTEMS =============================
\section{Optical Waveguides as Quantum Systems}

\subsection{Basics of Waveguide Coupling and Evanescent Fields}

Coupled optical waveguides represent a passive quantum platform where two single-mode waveguides interact through evanescent field overlap. The system is governed by coupled-mode theory \cite{laiNonclassicalFieldsLinear1991}.

The quantum Hamiltonian for coupled waveguides:
\begin{equation}
H = \hbar\omega(a^\dagger a + b^\dagger b) + \hbar J(a^\dagger b + b^\dagger a)
\end{equation}

where:
\begin{itemize}
	\item $J$ is the coupling strength (depends on waveguide separation)
	\item $a, b$ are annihilation operators for the two modes
\end{itemize}

\subsection{Losses in Realistic Systems}

Real waveguides experience material losses (absorption, scattering):
\begin{equation}
\gamma \text{ (loss rate)} = \frac{2.3 \times \alpha}{10} \text{ (in natural units from dB/cm)}
\end{equation}

Typical parameters:
\begin{itemize}
	\item Silica: $\gamma \approx 3 \times 10^9$ s$^{-1}$, $\gamma/J \approx 1/50$ (excellent)
	\item LiNbO$_3$: $\gamma \approx 3 \times 10^9$ s$^{-1}$, $\gamma/J \approx 1/7$ (moderate)
	\item AlGaAs: $\gamma \approx 2.7 \times 10^{10}$ s$^{-1}$, $\gamma/J \approx 1/10$ (moderate)
\end{itemize}

\subsection{Relevance for Quantum Circuits and Quantum Photonics}

Coupled waveguides serve as basic building blocks for integrated quantum circuits \cite{PolitiSilicaonSilicon2008}. Success in generating robust entanglement in these passive structures motivates investigation of steering—potentially a more useful resource for practical quantum information protocols.

% ============================= SECTION 2.7: ACTIVE QUANTUM SYSTEMS AND STEERING =============================
\section{Application and Motivation: Why Steering in Lossy Waveguides Matters}

\subsection{Role in Photonic Quantum Networks}

Quantum steering in optical systems enables:

\begin{itemize}
	\item \textbf{One-Sided Device-Independent (1sDI) Quantum Networks:} Asymmetric steering allows distributed quantum networks where only central nodes need trusted devices \cite{branciardOnesidedDeviceIndependentQuantum2012}.

	\item \textbf{Secure Communication:} The asymmetry of steering provides inherent directionality for secure quantum channels, particularly in untrusted or adversarial scenarios.

	\item \textbf{Quantum Internet and Cryptography:} Growing quantum internet frameworks (QIA, GSMA) identify steering-based protocols as key enablers for next-generation secure communications \cite{downlingQuantumTechnology2003}.

	\item \textbf{Quantum Key Distribution (QKD):} Device-independent QKD schemes benefit from steering's asymmetric nature, reducing hardware complexity.
\end{itemize}

\subsection{Steering vs Entanglement for Practical Applications}

Why focus on steering rather than entanglement?

\begin{itemize}
	\item \textbf{Resource Asymmetry:} Entanglement requires both parties to trust devices. Steering allows one untrusted party.

	\item \textbf{Operational Advantage:} Steering is a more practical resource for scenarios with asymmetric trust levels or capabilities.

	\item \textbf{Hierarchical Efficiency:} Since Steering $\subset$ Entanglement, steering-capable systems are automatically entanglement-capable, but the converse is not guaranteed.

	\item \textbf{Future-Oriented:} As quantum networks grow, asymmetric architectures (central trusted hubs, edge nodes) will dominate, making steering naturally suited for these topologies.
\end{itemize}

\subsection{Quantum Internet, Cryptography, and Beyond}

Emerging quantum communication infrastructure requires efficient, robust, and asymmetric quantum resources:

\begin{itemize}
	\item Quantum internet protocols prioritize asymmetric trust models
	\item Quantum cryptography standards (ETSI, ISO) increasingly focus on device-independent protocols
	\item Quantum sensing and metrology benefit from asymmetric correlations
	\item Quantum machine learning exploits steering for advanced protocols
\end{itemize}

% ============================= SECTION 2.8: CHAPTER SUMMARY =============================
\section{Chapter Summary}

This chapter has established the theoretical and motivational foundations for investigating quantum steering:

\begin{enumerate}
	\item \textbf{Quantum Correlations Hierarchy:} Discord $\supseteq$ Entanglement $\supseteq$ Steering $\supseteq$ Bell Nonlocality—steering occupies the optimal position for practical quantum information.

	\item \textbf{Continuous-Variable Advantage:} CV systems offer deterministic generation, optical compatibility, and Gaussian state formalism—ideal for steering studies.

	\item \textbf{Steering Fundamentals:} Asymmetric, directional, intermediate in strength, quantifiable for Gaussian states—uniquely suited for one-sided device-independent protocols.

	\item \textbf{Loss and Decoherence:} Prior work shows entanglement survives realistic loss in waveguides; steering robustness remains to be investigated.

	\item \textbf{Waveguides as Platforms:} Coupled waveguides demonstrate that passive systems can preserve entanglement. Active systems (like our RDQBL) may offer superior steering generation and control.

	\item \textbf{Practical Motivation:} Emerging quantum internet, cryptography, and sensing applications specifically require steering resources in asymmetric network topologies.
\end{enumerate}

The stage is now set for developing the specific RDQBL model (Chapter \ref{chp3}), investigating steering generation mechanisms (Chapter \ref{chp4}), and demonstrating practical feasibility (Chapter \ref{chp5}).



% ============================= CHAPTER 3: MODEL & METHODOLOGY =============================
%=============================================================================================
%                       Chapter 3: Physical Model and Theoretical Formulation
%=============================================================================================

\chapter{Raman-Driven Quantum Beat Laser: Model and Dynamics}
\label{chp3}
\newpage

\section{Introduction}

This chapter presents the detailed theoretical formulation of the Raman-Driven Quantum Beat Laser (RDQBL) system. We develop the complete quantum mechanical model, derive the master equation governing system evolution, and establish the mathematical framework for analyzing quantum steering in this active quantum system.

The chapter is organized as follows: first, we describe the physical system configuration and atomic level structure; next, we present the full Hamiltonian in the interaction picture; then, we derive the density matrix equations of motion including cavity losses and atomic decay; finally, we develop the covariance matrix formalism for quantifying quantum steering in the output field.

\section{Physical System Configuration}

\subsection{Four-Level Atomic System}

The RDQBL employs a four-level atomic system in cascade configuration with energy levels:
\begin{itemize}
	\item $|a\rangle$: Ground state (stable, lifetime $\rightarrow \infty$)
	\item $|b\rangle$: Intermediate state (stable, lifetime $\rightarrow \infty$)
	\item $|c\rangle$: Excited state (short-lived, decays rapidly)
\end{itemize}

The relevant dipole transitions are:
\begin{itemize}
	\item $|a\rangle \leftrightarrow |b\rangle$: Resonant with cavity mode 1 (frequency $\nu_1$)
	\item $|b\rangle \leftrightarrow |c\rangle$: Resonant with cavity mode 2 (frequency $\nu_2$)
	\item $|a\rangle \leftrightarrow |c\rangle$: Driven by external classical field (Raman process)
\end{itemize}

\subsection{Schematic Representation}

[Insert Figure 3.1: RDQBL System Schematic showing four-level configuration, cavity modes, and driving fields]

\section{The Hamiltonian}

\subsection{Full Hamiltonian in Interaction Picture}

In the electric dipole and rotating wave approximations, the interaction picture Hamiltonian is:
\begin{equation}
H = \hbar g_1 a_1 |a\rangle\langle b| + \hbar g_2 a_2 |b\rangle\langle c| - \frac{\hbar\Omega}{2}e^{-i\phi}|a\rangle\langle c| + \text{H.c.}
\end{equation}

where:
\begin{itemize}
	\item $a_1, a_2$: Annihilation operators for cavity modes 1 and 2
	\item $g_1, g_2$: Cavity-atom coupling strengths
	\item $\Omega e^{-i\phi}$: Rabi frequency and phase of external driving field
	\item H.c.: Hermitian conjugate
\end{itemize}

\subsection{Parameters and Their Physical Significance}

\textbf{Coupling Constants:} $g_1, g_2$ represent the strength of light-matter interaction. Typical values: $g_i \sim 10^6 - 10^7$ Hz.

\textbf{Rabi Frequency:} $\Omega$ characterizes the strength of the external driving field inducing the $|a\rangle \leftrightarrow |c\rangle$ transition. Controllable experimentally via laser intensity.

\textbf{Relative Phase:} $\phi$ is the relative phase between the two Raman driving fields. This parameter enables control over quantum steering directionality.

\section{Equations of Motion}

\subsection{Density Matrix Evolution}

The time evolution of density matrix elements is governed by the master equation including decay terms:
\begin{equation}
\dot{\rho}_{ij} = -\frac{i}{\hbar}\sum_k(V_{ik}\rho_{kj} - \rho_{ik}V_{kj}) - \frac{1}{2}\sum_k(\Gamma_{ik}\rho_{kj} + \rho_{ik}\Gamma_{kj})
\end{equation}

where $\Gamma_{ik}$ are decay coefficients for atomic transitions.

\subsection{Relevant Density Matrix Elements}

For the cascade system, the relevant coherences are:
\begin{itemize}
	\item $\rho_{ab} = \langle a|\rho|b\rangle$: Coherence between ground and intermediate states
	\item $\rho_{bc} = \langle b|\rho|c\rangle$: Coherence between intermediate and excited states
	\item $\rho_{ac} = \langle a|\rho|c\rangle$: Coherence between ground and excited states (driven by external field)
\end{itemize}

\subsubsection{Equations for Coherences}

\begin{align}
\dot{\rho}_{ab} &= -\gamma\rho_{ab} + \frac{i\Omega}{2}e^{-i\phi}\rho_{cb} - ig_1(a_1\rho_{bb} - \rho_{aa}a_1) + ig_2\rho_{ac}a_2^\dagger \\
\dot{\rho}_{bc} &= -\gamma\rho_{bc} - \frac{i\Omega}{2}e^{-i\phi}\rho_{ba} - ig_2(a_2\rho_{cc} - \rho_{bb}a_2) - ig_1a_1^\dagger\rho_{ac}
\end{align}

where $\gamma$ is the atomic decay rate for the excited state $|c\rangle$.

\section{Cavity Field Dynamics}

\subsection{Master Equation for Reduced Density Operator}

Tracing over atomic degrees of freedom and focusing on the cavity field, the evolution of $\rho_f$ (field reduced density matrix) is:
\begin{equation}
\dot{\rho}_f = \frac{1}{\hbar}[H_{\text{eff}}, \rho_f] + \mathcal{L}_{\text{cavity}}[\rho_f]
\end{equation}

where the cavity loss Liouvillian is:
\begin{equation}
\mathcal{L}_{\text{cavity}}[\rho_f] = \kappa_1(a_1\rho_f a_1^\dagger - \frac{1}{2}\{a_1^\dagger a_1, \rho_f\}) + \kappa_2(a_2\rho_f a_2^\dagger - \frac{1}{2}\{a_2^\dagger a_2, \rho_f\})
\end{equation}

with cavity decay rates $\kappa_1, \kappa_2$.

\section{Covariance Matrix Formulation}

\subsection{Quadrature Operators}

The quadrature operators for the two cavity modes are defined as:
\begin{align}
X_1 &= \frac{1}{\sqrt{2}}(a_1 + a_1^\dagger), \quad P_1 = \frac{i}{\sqrt{2}}(a_1^\dagger - a_1) \\
X_2 &= \frac{1}{\sqrt{2}}(a_2 + a_2^\dagger), \quad P_2 = \frac{i}{\sqrt{2}}(a_2^\dagger - a_2)
\end{align}

\subsection{Two-Mode Covariance Matrix}

For a two-mode Gaussian state, all correlations are captured by the covariance matrix:
\begin{equation}
V = \begin{pmatrix}
V_1 & C_{12} \\
C_{12}^T & V_2
\end{pmatrix}
\end{equation}

where:
\begin{align}
V_1 &= \begin{pmatrix}
\langle X_1^2 \rangle - \langle X_1 \rangle^2 & \frac{1}{2}(\langle X_1 P_1 + P_1 X_1 \rangle - 2\langle X_1\rangle\langle P_1 \rangle) \\
\frac{1}{2}(\langle X_1 P_1 + P_1 X_1 \rangle - 2\langle X_1\rangle\langle P_1 \rangle) & \langle P_1^2 \rangle - \langle P_1 \rangle^2
\end{pmatrix} \\
C_{12} &= \begin{pmatrix}
\langle X_1 X_2 \rangle - \langle X_1\rangle\langle X_2 \rangle & \frac{1}{2}(\langle X_1 P_2 + P_2 X_1 \rangle - 2\langle X_1\rangle\langle P_2 \rangle) \\
\frac{1}{2}(\langle X_2 P_1 + P_1 X_2 \rangle - 2\langle X_2\rangle\langle P_1 \rangle) & \langle P_1 P_2 \rangle - \langle P_1\rangle\langle P_2 \rangle
\end{pmatrix}
\end{align}

\section{Quantifying Quantum Steering}

\subsection{Steering Criterion for Gaussian States}

For the two-mode Gaussian state with covariance matrix $V$, steering from subsystem 1 to subsystem 2 is quantified as:
\begin{equation}
S^{1 \to 2} = \max\left\{0, \frac{1}{2}\log_2\frac{\det V_1}{4\det V_{\text{out}}^T}\right\}
\end{equation}

where $V_{\text{out}}^T$ is evaluated from the conditional state after projective measurements by party 1.

\subsection{Symplectic Eigenvalues}

For Gaussian states, the steering criterion can equivalently be expressed using symplectic eigenvalues of the partial transpose of the covariance matrix:
\begin{equation}
S^{1 \to 2} > 0 \iff \tilde{\nu}_- < \frac{1}{2}
\end{equation}

where $\tilde{\nu}_-$ is the smallest symplectic eigenvalue.

\subsection{Comparison with Entanglement}

Logarithmic negativity (entanglement):
\begin{equation}
E = \max\{0, -\log_2(2\tilde{\nu}_-)\}
\end{equation}

Key difference: Steering criterion depends on determinant of subsystem 1 covariance ($\det V_1$), while entanglement criterion depends on determinants of both subsystems and their cross-correlations.

\section{Initial Conditions and System Parameters}

\subsection{Input Cavity Field State}

We consider the initial cavity field as a product of two independent single-mode Gaussian states:
\begin{equation}
\rho_{\text{in}} = \rho_1 \otimes \rho_2
\end{equation}

Each single-mode state is characterized by:
\begin{itemize}
	\item Non-classicality parameter: $\tau$ (ranging 0 to 1/2 for squeezed states)
	\item Purity: $\mu$ (ranging 0 for maximally mixed to 1 for pure)
\end{itemize}

\subsection{Typical Parameters}

Realistic RDQBL parameters for simulations:
\begin{itemize}
	\item Coupling strengths: $g_1 = g_2 = 16$ kHz
	\item External field Rabi frequency: $\Omega = 2160$-$3160$ kHz (tunable)
	\item Relative phase: $\phi \in [0, 2\pi]$ (tunable)
	\item Atomic decay rate: $\gamma = 20$ kHz
	\item Cavity decay rates: $\kappa_1, \kappa_2 = 0$-$0.1$ kHz
	\item Atom injection rate: $r_a = 10$ kHz
\end{itemize}

\section{Chapter Summary}

This chapter has provided:
\begin{enumerate}
	\item Complete description of the RDQBL four-level system configuration
	\item Full quantum Hamiltonian in interaction picture
	\item Master equations for cavity field evolution including losses
	\item Covariance matrix formalism for Gaussian state analysis
	\item Mathematical criteria for steering quantification
	\item Physical interpretation of tunable parameters
	\item Realistic system parameters for numerical studies
\end{enumerate}

These foundations enable the numerical investigations presented in Chapter \ref{chp4}, where we solve the system dynamics and analyze steering generation and control.



% ============================= CHAPTER 4: RESULTS & ANALYSIS =============================
%=============================================================================================
%                               Chapter 4: Results and Analysis
%=============================================================================================

\chapter{Quantum Steering Generation and Dynamics}
\label{chp4}
\newpage

\section{Introduction}

This chapter presents comprehensive numerical and analytical results investigating continuous-variable quantum steering in the RDQBL system. We systematically explore how steering is generated, sustained, and controlled through various system parameters. The results are organized to address the key research questions: (1) Can steering be generated reliably? (2) What parameters control steering directionality? (3) How robust is steering under realistic dissipation?

\section{Time Evolution of Quantum Steering}

\subsection{Baseline Scenario: Lossless System}

Figure \ref{fig:steering_evolution} presents the time evolution of steering signals $S^{1 \to 2}$ and $S^{2 \to 1}$ for the lossless case.

[Insert Figure 4.1: Time evolution of steering $S^{1 \to 2}$ and $S^{2 \to 1}$ vs dimensionless time $\Omega t$]

\textbf{Key Observations:}
\begin{itemize}
	\item Steering emerges and reaches maximum around $\Omega t \approx \pi/4$
	\item Clear asymmetry: $S^{1 \to 2} > S^{2 \to 1}$ for certain time windows
	\item Oscillatory behavior reflects quantum beat dynamics
	\item Both steering signals vanish periodically as predicted by theoretical model
\end{itemize}

\section{Effect of Rabi Frequency}

\subsection{Variation with External Driving Field}

The Rabi frequency $\Omega$ of the external Raman driving field is a critical parameter.

[Insert Figure 4.2: Steering vs time for varying $\Omega$: 2160 kHz, 2660 kHz, 3160 kHz]

\textbf{Findings:}
\begin{itemize}
	\item Increased $\Omega$ enhances steering strength (higher $S_{\max}$)
	\item Time duration of steering window increases with $\Omega$
	\item Maximum steering is achieved at $\Omega \approx 2660$ kHz for these parameters
	\item Trade-off: Very large $\Omega$ reduces steering time window
\end{itemize}

\section{Effect of Relative Phase Control}

\subsection{Phase-Dependent Steering Directionality}

The relative phase $\phi$ between the two Raman driving fields provides direct control over steering asymmetry.

[Insert Figure 4.3: Phase-dependent steering: (a) $\phi = 0$, (b) $\phi = \pi/4$, (c) $\phi = \pi/2$]

\textbf{Key Results:}
\begin{itemize}
	\item At $\phi = 0$: Nearly symmetric steering $S^{1 \to 2} \approx S^{2 \to 1}$
	\item At $\phi = \pi/2$: Strong asymmetry, $S^{1 \to 2} \gg S^{2 \to 1}$
	\item Steering directionality is continuously tunable through $\phi$
	\item Optimal asymmetry parameter regimes can be identified
\end{itemize}

\section{Cavity Damping Effects}

\subsection{Loss-Induced Steering Degradation}

Figure \ref{fig:damping} shows how cavity damping ($\kappa_1, \kappa_2$) affects steering survival.

[Insert Figure 4.4: Steering decay with increasing cavity damping rates]

\textbf{Observations:}
\begin{itemize}
	\item Moderate damping ($\kappa \lesssim 0.005$ kHz) shows minimal steering degradation
	\item At $\kappa = 0.01$ kHz, steering magnitude reduced by $\sim 30\%$
	\item Steering persists even at experimentally realistic damping rates
	\item Time-window for steering detection shrinks with increased loss
\end{itemize}

\subsection{Comparison: Steering vs Entanglement under Loss}

[Insert Figure 4.5: Comparative robustness of steering vs entanglement]

\textbf{Finding:} Steering shows comparable or superior robustness compared to entanglement under cavity damping, supporting the thesis that active systems effectively preserve quantum correlations.

\section{Non-Classicality and Purity Effects}

\subsection{Initial State Non-Classicality}

The non-classicality parameter $\tau$ of initial cavity modes significantly influences steering dynamics.

[Insert Figure 4.6: Effect of non-classicality on steering: $\tau = 0.34, 0.37, 0.40$]

\textbf{Results:}
\begin{itemize}
	\item Increased $\tau$ (higher non-classicality) enhances steering strength
	\item Non-classical initial states generate steering more efficiently
	\item Steering time-window extends with increased non-classicality
	\item Effect is more pronounced for steering than for entanglement
\end{itemize}

\subsection{Purity of Initial States}

[Insert Figure 4.7: Effect of purity on steering: $\mu = 0.75, 0.85, 1.0$]

\textbf{Key Findings:}
\begin{itemize}
	\item Pure initial states ($\mu = 1$) show maximum steering
	\item Mixed initial states ($\mu = 0.75$) still generate steering but at reduced magnitude
	\item Purity effect is more pronounced than non-classicality effect for steering
	\item Practical implication: High-purity initial states required for robust steering
\end{itemize}

\section{Hierarchy Verification: Discord $\supseteq$ Entanglement $\supseteq$ Steering}

\subsection{Temporal Ordering of Quantum Correlations}

[Insert Figure 4.8: Time evolution of all four quantum correlations: Discord, Entanglement, Steering, Bell Nonlocality]

\textbf{Verification:}
\begin{itemize}
	\item Discord exhibits longest survival time, reaching steady state
	\item Entanglement emerges after discord and persists longer than steering
	\item Steering appears after entanglement and vanishes earlier
	\item Bell nonlocality (if present) shows shortest lifetime
	\item Temporal order confirms theoretical hierarchy
\end{itemize}

\subsection{Boundary Regions}

[Insert Table 4.1: Parameter ranges for correlation existence]

Clear identification of parameter regimes where only certain correlations exist supports the strict hierarchy.

\section{Optimized Steering Regimes}

\subsection{Parameter Space Mapping}

[Insert Figure 4.9: 2D parameter map showing steering strength as function of ($\Omega$, $\phi$)]

This contour plot identifies optimal parameter combinations for maximum steering with controlled asymmetry.

\textbf{Practical Guidelines:}
\begin{itemize}
	\item For symmetric steering: $\phi \approx 0$, $\Omega \approx 2400$ kHz
	\item For asymmetric steering (Mode 1 steering Mode 2): $\phi \approx \pi/2$, $\Omega \approx 2660$ kHz
	\item For reverse asymmetry: $\phi \approx 3\pi/2$, $\Omega \approx 2660$ kHz
\end{itemize}

\section{Comparison with CEL Results}

\subsection{RDQBL vs Correlated Emission Laser}

[Insert Figure 4.10: Steering comparison between RDQBL (this work) and CEL (Ullah et al. 2019)]

\textbf{Advantages of RDQBL:}
\begin{itemize}
	\item Stronger directional control through phase tuning
	\item Enhanced robustness against cavity damping
	\item More efficient steering generation per unit Rabi frequency
	\item Better compatibility with quantum beat control
\end{itemize}

\textbf{Trade-offs:}
\begin{itemize}
	\item RDQBL requires phase stabilization (more stringent experimental requirement)
	\item CEL may be simpler to implement in some configurations
\end{itemize}

\section{Analytical Validation}

\subsection{Comparison with Strongly-Driven Limit Approximation}

In the strongly driven limit ($\Omega \gg \gamma$), we derived analytical expressions for key observables.

[Insert Figure 4.11: Analytical vs Numerical solutions for steering evolution]

\textbf{Validation Result:}
Maximum deviation between analytical approximation and full numerical solution: $\sim 5\%$ for $\Omega \geq 2160$ kHz, confirming validity of approximation scheme.

\section{Chapter Summary}

Key findings of this chapter:

\begin{enumerate}
	\item Continuous-variable quantum steering can be reliably generated in RDQBL systems
	\item Steering strength and directionality are tunable through accessible parameters ($\Omega$, $\phi$)
	\item Relative phase provides unprecedented control over steering asymmetry
	\item Steering shows robustness comparable to or exceeding entanglement under realistic losses
	\item Quantum correlation hierarchy (Discord $\supseteq$ Entanglement $\supseteq$ Steering) is verified numerically
	\item Optimal parameter regimes identified for practical implementation
	\item RDQBL offers advantages over existing platforms for steering generation
\end{enumerate}

These results establish the RDQBL as a promising platform for steering-based quantum technologies and provide a roadmap for experimental realization, discussed further in Chapter \ref{chp5}.



% ============================= CHAPTER 5: DISCUSSION =============================
%=============================================================================================
%                                 Chapter 5: Discussion
%=============================================================================================

\chapter{Discussion and Physical Interpretation}
\label{chp5}
\newpage

\section{Introduction}

The numerical results presented in Chapter \ref{chp4} demonstrate the successful generation and control of continuous-variable quantum steering in RDQBL systems. This chapter provides deeper physical interpretation of these findings, discusses their implications for quantum technologies, and contextualizes them within the broader landscape of quantum information science.

\section{Mechanisms of Steering Generation in RDQBL}

\subsection{Role of Quantum Beat Interference}

The quantum beat frequency $\Omega = \omega_1 - \omega_2$ plays a central role in steering generation. The periodic modulation of cavity field correlations at the beat frequency enables the stepwise build-up of asymmetric steering.

\textbf{Physical Mechanism:}
\begin{enumerate}
	\item Two Raman driving fields create coherent superpositions in the atomic population
	\item Quantum beats emerge in the amplitude of Raman transitions
	\item These beats modulate the effective coupling to the two cavity modes
	\item Differential modulation between modes 1 and 2 creates asymmetric correlations
	\item Result: Directional steering from one mode to the other
\end{enumerate}

\subsection{Phase Control of Steering Directionality}

The relative phase $\phi$ between Raman fields provides unprecedented control:

\textbf{Theoretical Explanation:}
The Hamiltonian contains terms proportional to $e^{-i\phi}$ and $e^{i\phi}$, which directly control the interference pattern between Raman transitions. By tuning $\phi$, one can selectively enhance or suppress specific interference pathways, thereby steering quantum correlations preferentially from one mode to the other.

\textbf{Practical Consequence:}
Complete directionality control enables experimental switching between $S^{1 \to 2}$-dominant and $S^{2 \to 1}$-dominant regimes continuously and reversibly.

\section{Robustness Analysis Under Realistic Conditions}

\subsection{Why RDQBL Steering Survives Cavity Damping}

The robustness of RDQBL steering against cavity losses (Section \ref{sec:cavdamp}) can be understood through several factors:

\subsubsection{Coherence Preservation via Raman Process}

Unlike direct optical transitions, Raman processes bypass the excited state $|c\rangle$, which has the shortest lifetime. Population transfer occurs between long-lived states $|a\rangle$ and $|b\rangle$, creating robust coherence that:
\begin{itemize}
	\item Persists despite cavity photon loss
	\item Continuously regenerates correlated photons
	\item Maintains quantum phase relationships needed for steering
\end{itemize}

\subsubsection{Active System Advantage}

The RDQBL actively generates steering through coherent driving fields, in contrast to passive systems that merely redistribute pre-existing correlations:
\begin{itemize}
	\item Active systems continuously compensate for dissipation
	\item Gain processes (stimulated emission) can overcome losses
	\item System operates in non-equilibrium steady state with energy input
\end{itemize}

\subsection{Critical Loss Thresholds}

Analysis of steering decay rates reveals critical thresholds:

\textbf{Steering Survival Criterion:}
\begin{equation}
\frac{\kappa}{\Omega} \lesssim 0.005
\end{equation}

Below this ratio, steering shows only modest degradation. This corresponds to cavity quality factors:
\begin{equation}
Q_{\text{cavity}} = \frac{\Omega}{\kappa} \gtrsim 200
\end{equation}

\textbf{Experimental Perspective:} High-finesse optical cavities routinely achieve $Q > 10^4$, providing ample margin for steering generation and detection.

\section{Hierarchy Implications and Quantum Correlations Ontology}

\subsection{Why Steering Occupies a Unique Position}

The strict hierarchy Discord $\supseteq$ Entanglement $\supseteq$ Steering $\supseteq$ Bell Nonlocality reflects fundamental quantum mechanical structures:

\textbf{Set-Theoretic Interpretation:}
\begin{itemize}
	\item Every steerable state is entangled (Steering $\subset$ Entanglement)
	\item Every entangled state has discord (Entanglement $\subset$ Discord)
	\item But the reverse implications don't hold: separable states can exhibit discord; entangled states may not be steerable
\end{itemize}

\textbf{Why Steering for Quantum Technology?}
Steering combines advantages from both stronger and weaker correlations:
\begin{itemize}
	\item Stronger than Bell nonlocality: More prevalent, easier to generate
	\item Weaker than entanglement: More robust, survives under certain conditions
	\item Asymmetric nature: Uniquely suited for asymmetric trust models in quantum networks
	\item Intermediate strength: Provides sweet spot for practical applications
\end{itemize}

\section{Implications for Quantum Information Processing}

\subsection{One-Sided Device-Independent (1sDI) Quantum Key Distribution}

Steering's asymmetric nature is essential for 1sDI-QKD protocols. In such schemes:

\textbf{Setup:}
\begin{itemize}
	\item Alice operates a trusted device (prepared states, reliable measurements)
	\item Bob operates an untrusted device (or is even adversarial)
	\item Security relies on Bob's inability to prepare correlations that explain Alice's measurement results
\end{itemize}

\textbf{Security Analysis:}
Security is guaranteed if Alice can steer Bob's conditional state in a manner incompatible with local hidden variables. RDQBL steering directly provides this resource.

\textbf{Key Rate Calculation (Simplified):}
\begin{equation}
R_{\text{1sDI-QKD}} \propto S^{\text{Alice} \to \text{Bob}}
\end{equation}

Stronger steering directly translates to higher secure key rates.

\subsection{Asymmetric Quantum Networks and Photonic Architectures}

As quantum internet infrastructure develops, network topologies will likely be asymmetric:
\begin{itemize}
	\item Central trusted hubs with guaranteed device quality
	\item Peripheral nodes with reduced resources or variable trust levels
	\item Network links connecting heterogeneous nodes
\end{itemize}

RDQBL-based steering sources enable efficient operation in such architectures.

\section{Comparison with Prior Work}

\subsection{Waveguides: Entanglement in Passive Systems}

Previous work (Rai et al., Optics Express 2010) demonstrated entanglement survival in coupled lossy waveguides. Key comparison:

\begin{table}[h]
\centering
\begin{tabular}{lcc}
\textbf{Property} & \textbf{Waveguides} & \textbf{RDQBL} \\
\hline
System Type & Passive & Active \\
Entanglement Generation & Requires parametric source & Intrinsic to system \\
Tuning Capability & Limited & Extensive ($\Omega$, $\phi$, etc.) \\
Directionality & Not applicable & Controllable \\
Robustness & Good & Excellent \\
\end{tabular}
\end{table}

\textbf{Conclusion:} While waveguides preserve entanglement, RDQBL systems actively generate tunable steering—a stronger and more practically useful resource.

\subsection{CEL: Steering in Active Systems}

Recent work (Ullah et al., Optics Express 2019) established steering generation in Correlated Emission Lasers. Our RDQBL work extends this:

\begin{table}[h]
\centering
\begin{tabular}{lcc}
\textbf{Property} & \textbf{CEL} & \textbf{RDQBL} \\
\hline
Atomic Configuration & Three-level cascade & Four-level cascade + Raman \\
Driving Mechanism & Direct transitions & Raman process \\
Directionality Control & Rabi frequency & Rabi frequency + Phase \\
Robustness & Moderate & Enhanced \\
Coherence Preservation & Via cascade & Via long-lived states + Raman \\
\end{tabular}
\end{table}

\textbf{Innovation:} RDQBL adds phase-based steering control, enabling unprecedented directional selectivity.

\section{Experimental Feasibility and Implementation Roadmap}

\subsection{Realizable Physical Platforms}

Several experimental platforms can implement RDQBL:

\subsubsection{Alkali Atoms (Rb, Cs)}
\begin{itemize}
	\item Pros: Well-understood level structures, long coherence times, existing experimental expertise
	\item Cons: Requires high-finesse cavities, temperature stabilization
	\item Feasibility: High (multiple groups worldwide have demonstrated similar systems)
\end{itemize}

\subsubsection{Trapped Ions}
\begin{itemize}
	\item Pros: Excellent coherence, perfect quantum control
	\item Cons: Requires specialized equipment, scalability challenges
	\item Feasibility: Moderate (requires dedicated ion-trap facility)
\end{itemize}

\subsubsection{Solid-State Defect Centers (NV in Diamond)}
\begin{itemize}
	\item Pros: Room-temperature operation, compact systems
	\item Cons: Shorter coherence times, complex level structures
	\item Feasibility: Moderate (emerging technology)
\end{itemize}

\subsection{Experimental Challenges and Mitigation Strategies}

\textbf{Challenge 1: Phase Stability}
\begin{itemize}
	\item Issue: Relative phase $\phi$ must be maintained with sub-radian precision
	\item Solution: Active stabilization using feedback locks (e.g., Pound-Drever-Hall)
	\item Precedent: Standard in modern quantum optics labs
\end{itemize}

\textbf{Challenge 2: Cavity Loss}
\begin{itemize}
	\item Issue: Steering requires $Q > 200$
	\item Solution: Ultra-high-finesse cavities (Q $> 10^6$ readily available)
	\item Precedent: Standard in cavity QED experiments
\end{itemize}

\textbf{Challenge 3: Temperature and Vibration Stability}
\begin{itemize}
	\item Issue: Thermal drifts shift transition frequencies
	\item Solution: Precision stabilization systems (available commercially)
	\item Precedent: Implemented in cold-atom labs worldwide
\end{itemize}

\subsection{Step-by-Step Implementation Path}

\begin{enumerate}
	\item \textbf{Year 1}: Demonstrate steering generation in lossless regime using alkali atoms
	\item \textbf{Year 2}: Add cavity losses and verify robustness predictions
	\item \textbf{Year 3}: Demonstrate phase-controlled directionality switching
	\item \textbf{Year 4}: Integrate with quantum cryptography protocols for proof-of-principle QKD
	\item \textbf{Year 5}: Optimize for integration into quantum network nodes
\end{enumerate}

\section{Limitations and Future Directions}

\subsection{Current Limitations}

\begin{enumerate}
	\item \textbf{Single-Mode Analysis:} Current work assumes perfect cavity modes; real cavities have spatial variations
	\item \textbf{Gaussian Assumption:} Non-Gaussian corrections not analyzed
	\item \textbf{No Measurement Details:} Actual homodyne detection limitations not modeled
	\item \textbf{Atom Number Approximation:} Assumes many atoms; single-atom effects not explored
\end{enumerate}

\subsection{Extensions for Future Work}

\begin{enumerate}
	\item \textbf{Multi-Mode Systems:} Extend to three or more cavity modes for quantum networks
	\item \textbf{Non-Gaussian Effects:} Include higher-order corrections and non-Gaussian operations
	\item \textbf{Measurement Device Noise:} Model realistic homodyne detector imperfections
	\item \textbf{Atom Dynamics:} Investigate collective effects and atom number fluctuations
	\item \textbf{Quantum Repeaters:} Integrate RDQBL with quantum repeater protocols
\end{enumerate}

\section{Chapter Summary}

This chapter has provided:
\begin{enumerate}
	\item Physical mechanisms explaining steering generation in RDQBL
	\item Analysis of robustness under realistic dissipation conditions
	\item Interpretation of quantum correlation hierarchy through set theory
	\item Applications to quantum cryptography and networks
	\item Comparison with prior work on entanglement and steering
	\item Experimental feasibility assessment and implementation roadmap
	\item Discussion of limitations and future research directions
\end{enumerate}

These discussions establish RDQBL not just as a theoretical model, but as a practically realizable platform for steering-based quantum technologies.



% ============================= CHAPTER 6: CONCLUSIONS =============================
%=============================================================================================
%                                 Chapter 6: Conclusions
%=============================================================================================

\chapter{Conclusions and Future Perspectives}
\label{chp6}
\newpage

\section{Summary of Main Findings}

This dissertation has investigated continuous-variable quantum steering in Raman-driven quantum beat laser (RDQBL) systems. The main contributions are:

\begin{enumerate}
	\item Developed comprehensive theoretical treatment of RDQBL system with steering quantification framework
	\item Demonstrated phase-controlled tunability of steering directionality
	\item Verified quantum correlation hierarchy in active laser systems
	\item Established robustness of steering under realistic cavity damping
\end{enumerate}

\section{Key Results}

\begin{itemize}
	\item Continuous-variable quantum steering can be reliably generated in RDQBL
	\item Steering directionality is tunable through relative phase control
	\item RDQBL steering shows superior robustness compared to waveguide systems
	\item Quantum correlation hierarchy Discord $\supseteq$ Entanglement $\supseteq$ Steering $\supseteq$ Bell Nonlocality is verified
\end{itemize}

\section{Significance and Impact}

This work provides:
\begin{enumerate}
	\item First comprehensive treatment of steering in Raman-driven systems
	\item Practical pathway for steering-based quantum technologies
	\item Roadmap for experimental implementation
	\item Foundation for asymmetric quantum networks
\end{enumerate}

\section{Limitations}

\begin{enumerate}
	\item Theoretical model assumes idealized conditions
	\item Gaussian state assumption (non-Gaussian corrections not included)
	\item Two-mode focus (extensions to multi-mode systems needed)
	\item Experimental validation pending
\end{enumerate}

\section{Recommendations for Future Work}

\subsection{Immediate Extensions}

\begin{itemize}
	\item Experimental demonstration of RDQBL steering
	\item Integration with quantum cryptography protocols
	\item Non-Gaussian analysis of steering dynamics
	\item Multi-mode cavity system extension
\end{itemize}

\subsection{Medium-Term Research}

\begin{itemize}
	\item Quantum repeater integration
	\item Steering-enhanced quantum sensing
	\item Implementation on alternative platforms (trapped ions, NV centers)
	\item Device-independent protocol demonstrations
\end{itemize}

\subsection{Long-Term Vision}

\begin{itemize}
	\item Deployment in quantum internet nodes
	\item Scalable distributed quantum computing with steering resources
	\item Fundamental tests of quantum mechanics
	\item Hybrid quantum systems combining multiple platforms
\end{itemize}

\section{Closing Remarks}

RDQBL systems provide a natural, tunable, and robust platform for generating continuous-variable quantum steering—a key resource for next-generation quantum technologies. This dissertation establishes the theoretical foundation; experimental confirmation and technological implementation await future researchers.

The quantum internet is emerging, and steering-based quantum resources will play a crucial role in its architecture.



% ============================= REFERENCES =============================
%=============================================================================================
%                                     Chapter 7: References
%=============================================================================================

\chapter*{References}
\addcontentsline{toc}{chapter}{References}

\begin{thebibliography}{99}

\bibitem{SchrodingerE35}
E. Schrödinger, ``Discussion of probability relations between separated systems,'' \textit{Mathematical Proceedings of the Cambridge Philosophical Society}, vol. 31, no. 4, pp. 555--563, 1935.

\bibitem{einsteinPodolskyRosen1935}
A. Einstein, B. Podolsky, and N. Rosen, ``Can quantum-mechanical description of physical reality be considered complete?,'' \textit{Physical Review}, vol. 47, no. 10, pp. 777--780, 1935.

\bibitem{BranciardC12}
C. Branciard, E. G. Cavalcanti, S. P. Walborn, V. Scarani, and H. M. Wiseman, ``One-sided device-independent quantum key distribution: Security, feasibility, and the connection with steering,'' \textit{Physical Review A}, vol. 85, no. 1, p. 010301(R), 2012.

\bibitem{WisemanH07}
H. M. Wiseman, S. J. Jones, and A. C. Doherty, ``Steering, entanglement, nonlocality, and the Einstein-Podolsky-Rosen paradox,'' \textit{Physical Review Letters}, vol. 98, no. 14, p. 140402, 2007.

\bibitem{WeedbrookC12}
C. Weedbrook, S. Pirandola, R. García-Patrón, N. J. Cerf, T. C. Ralph, J. H. Shapiro, and S. Lloyd, ``Gaussian quantum information,'' \textit{Reviews of Modern Physics}, vol. 84, no. 2, p. 621, 2012.

\bibitem{Qamar2008}
S. Qamar, F. Ghafoor, M. Hillery, and M. S. Zubairy, ``Quantum beat laser as a source of entangled radiation,'' \textit{Physical Review A}, vol. 77, no. 6, p. 062308, 2008.

\bibitem{Scully1997}
M. O. Scully and M. S. Zubairy, \textit{Quantum Optics}. Cambridge: Cambridge University Press, 1997.

\bibitem{Qamar2009}
S. Qamar, S. Al-Amri, and M. S. Zubairy, ``Entangled radiation via a Raman-driven quantum-beat laser,'' \textit{Physical Review A}, vol. 80, no. 3, p. 033818, 2009.

\bibitem{Zubairy1992}
M. S. Zubairy, M. O. Scully, and J. H. Eberly, ``Directional spontaneous emission from an extended ensemble of N atoms: Cooperation and disorder,'' \textit{Physical Review A}, vol. 36, no. 6, p. 2519, 1992.

\bibitem{ShakirUllah19}
S. Ullah, H. S. Qureshi, and F. Ghafoor, ``Coherence control of entanglement dynamics of two-mode Gaussian state via Raman driven quantum beat laser using Simon's criterion,'' \textit{Applied Optics}, vol. 58, no. 2, pp. 197--204, 2019.

\bibitem{RaiDasAgarwal2010}
A. Rai, S. Das, and G. S. Agarwal, ``Quantum entanglement in coupled lossy waveguides,'' \textit{Optics Express}, vol. 18, no. 6, pp. 6241--6254, 2010.

\bibitem{UllahQureshiGhafoor2019}
S. Ullah, H. S. Qureshi, and F. Ghafoor, ``Hierarchy of temporal quantum correlations using a correlated spontaneous emission laser,'' \textit{Optics Express}, vol. 27, no. 19, pp. 26858--26873, 2019.

\bibitem{kogiastQuantificationGaussianQuantum2015}
I. Kogias, A. R. Lee, S. Ragy, and G. Adesso, ``Quantification of Gaussian quantum steering,'' \textit{Physical Review Letters}, vol. 114, no. 6, p. 060403, 2015.

\bibitem{brunnerBellNonlocality2014}
N. Brunner, D. Cavalcanti, S. Pironio, V. Scarani, and S. Wehner, ``Bell nonlocality,'' \textit{Reviews of Modern Physics}, vol. 86, no. 2, p. 419, 2014.

\bibitem{PolitiSilicaonSilicon2008}
A. Politi, M. J. Cryan, J. G. Rarity, S. Yu, and J. L. O'Brien, ``Silica-on-silicon waveguide quantum circuits,'' \textit{Science}, vol. 320, no. 5876, pp. 646--649, 2008.

\bibitem{laiNonclassicalFieldsLinear1991}
W. K. Lai, V. Bužek, and P. L. Knight, ``Nonclassical fields in a linear directional coupler,'' \textit{Physical Review A}, vol. 43, no. 12, p. 6323, 1991.

\bibitem{downlingQuantumTechnology2003}
J. P. Dowling and G. J. Milburn, ``Quantum technology: the second quantum revolution,'' \textit{Philosophical Transactions of the Royal Society A: Mathematical, Physical and Engineering Sciences}, vol. 361, no. 1809, pp. 1655--1674, 2003.

\bibitem{VidalWerner2002}
G. Vidal and R. F. Werner, ``Computable measure of entanglement,'' \textit{Physical Review A}, vol. 65, no. 3, p. 032314, 2002.

\bibitem{AdessoDatta2010}
G. Adesso and A. Datta, ``Quantum versus classical correlations in Gaussian states,'' \textit{Physical Review Letters}, vol. 105, no. 3, p. 030501, 2010.

\bibitem{GiordaParis2010}
P. Giorda and M. G. A. Paris, ``Gaussian quantum discord,'' \textit{Physical Review Letters}, vol. 105, no. 20, p. 020503, 2010.

\bibitem{OlivierZurek2001}
H. Ollivier and W. H. Zurek, ``Quantum discord: a measure of the quantumness of correlations,'' \textit{Physical Review Letters}, vol. 88, no. 1, p. 017901, 2001.

\bibitem{HorodockiHorodocki2009}
R. Horodecki, P. Horodecki, M. Horodecki, and K. Horodecki, ``Quantum entanglement,'' \textit{Reviews of Modern Physics}, vol. 81, no. 2, p. 865, 2009.

\bibitem{ModiBrodutch2012}
K. Modi, A. Brodutch, H. Cable, T. Paterek, and V. Vedral, ``The classical-quantum boundary for correlations: Discord and related measures,'' \textit{Reviews of Modern Physics}, vol. 84, no. 4, p. 1655, 2012.

\bibitem{DuanGiedke2000}
L.-M. Duan, G. Giedke, J. I. Cirac, and P. Zoller, ``Inseparability criterion for continuous variable systems,'' \textit{Physical Review Letters}, vol. 84, no. 12, p. 2722, 2000.

\bibitem{SimonPeres2000}
R. Simon, ``Peres-Horodecki separability criterion for continuous variable systems,'' \textit{Physical Review Letters}, vol. 84, no. 12, p. 2726, 2000.

\bibitem{BarnettRadmore2002}
S. M. Barnett and P. M. Radmore, \textit{Methods in Theoretical Quantum Optics}. Oxford: Oxford University Press, 2002.

\bibitem{BraunsteinVanLoock2005}
S. L. Braunstein and P. van Loock, ``Quantum information with continuous variables,'' \textit{Reviews of Modern Physics}, vol. 77, no. 2, p. 513, 2005.

\end{thebibliography}



% ============================= APPENDICES (OPTIONAL) =============================
% Uncomment if needed
% \appendix
% \input{./Appendix_A}

\end{document}
